%% kapa: exam=2 round=24 year=2013 period=1 subject=appraisal_practice subject_ko=감정평가실무 seq=4 number=4 points=15
%% source: Q-Net 감정평가사 기출문제 (cst00309) artlSeq=5077404 file=extracted/1008/5077404/2013년 제24회 감정평가사 2차시험 기출문제/1교시.hwp
%% conversion: kordoc 4.24.51 -> tools/gen2.py (auto); review=proofread-sonnet; flags=none
\begin{longquestion}{【문제 4】}
감정평가사 홍길동 씨는 법원으로부터 시장가치 및 예상낙찰가 산정을 의뢰 받았다. 주어진 자료를 참고하여 시장가치 및 예상낙찰가를 구하고, 산출방법에 대해 약술하시오. (15점) \par
\textbf{<자료 1> 평가대상 부동산 개황} \par
1. 소재지: A시 B구 C동 100-1번지 외 \par
2. 토지: 100-1번지, 200m\textsuperscript{2}, 대, 준주거지역, 정방형, 평지 \par
100-2번지, 200m\textsuperscript{2}, 대, 준주거지역, 정방형, 평지 \par
100-3번지, 200m\textsuperscript{2}, 대, 준주거지역, 정방형, 평지 \par
3. 건물: 평가대상 토지(100-1, 2, 3번지) 및 평가대상 외 토지(100번지: 타인소유) 4필지 일단의 토지 지상 철근콘크리트조 슬래브지붕 1～4층 각 480m\textsuperscript{2}(지하층 없음) 근린생활시설(신축년도 2003. 9. 7) \par
4. 기준시점: 2013. 9. 7 \par
\textbf{<자료 2> 주위환경 및 시장상황} \par
1. 본건이 속한 A시 B구 C동은 전면 도로변으로 4～6층 근린생활시설이 혼재하며, 후면은 근린생활시설 및 단독, 다가구 등 주상복합지대로 형성되어 있음 \par
2. 본건 주위는 최근 2년간 가격변동추이는 보합정도이며, 향후 전망도 보합정도이나 용적률이 낮은 오래된 상업용 건물(1970년대 신축)의 경우 개별 또는 합필하여 철거 후 재건축이 진행 중인 필지도 일부 혼재함 \par
3. 토지소유자와 건물소유자가 다른 경우 건물소유자는 토지의 시장가치에 적정지료를 지불하고 정상적으로 사용ㆍ수익할 수 있는 것으로 조사되었으며, 주변의 토지 거래량이나 거래가격도 적정한 것으로 조사됨 \par
\textbf{<자료 3> 지적도} \par
\qfigure{figures/04-1.png}
\textbf{<자료 4> 표준지 공시지가} \par
(2013. 1. 1) \par
{\small\centering
\begin{tabular}{|c|c|c|c|c|c|c|c|c|}
\hline
\makecell{기호} & \makecell{소재지} & \makecell{면적\newline{}(m\textsuperscript{2})} & \makecell{지목} & \makecell{용도\newline{}지역} & \makecell{이용\newline{}상황} & \makecell{도로\newline{}교통} & \makecell{형상\newline{}지세} & \makecell{공시지가(원/m\textsuperscript{2})} \\
\hline
\makecell{1} & \makecell{C동 99-1} & \makecell{200} & \makecell{대} & \makecell{준주거} & \makecell{상업용} & \makecell{중로\newline{}한면} & \makecell{정방형\newline{}평지} & \makecell{5,200,000} \\
\hline
\makecell{2} & \makecell{C동 101-1} & \makecell{250} & \makecell{대} & \makecell{준주거} & \makecell{주상용} & \makecell{소로\newline{}한면} & \makecell{가장형\newline{}평지} & \makecell{3,000,000} \\
\hline
\makecell{3} & \makecell{C동 103} & \makecell{400} & \makecell{대} & \makecell{준주거} & \makecell{상업용} & \makecell{중로\newline{}한면} & \makecell{세장형\newline{}완경사} & \makecell{4,000,000} \\
\hline
\end{tabular}\par}

\textbf{<자료 5> 개별요인 비교치} \par
1. 접면도로 \par
{\small\centering
\begin{tabular}{|c|c|c|c|c|}
\hline
\makecell{구분} & \makecell{중로각지} & \makecell{중로한면} & \makecell{소로각지} & \makecell{소로한면} \\
\hline
\makecell{비교치} & \makecell{1.05} & \makecell{1.00} & \makecell{0.90} & \makecell{0.85} \\
\hline
\end{tabular}\par}

2. 형상 \par
{\small\centering
\begin{tabular}{|c|c|c|c|c|}
\hline
\makecell{구분} & \makecell{가로장방형} & \makecell{정방형} & \makecell{세로장방형} & \makecell{부정형} \\
\hline
\makecell{비교치} & \makecell{1.05} & \makecell{1.00} & \makecell{0.95} & \makecell{0.85} \\
\hline
\end{tabular}\par}

3. 지세 \par
{\small\centering
\begin{tabular}{|c|c|c|c|c|c|}
\hline
\makecell{구분} & \makecell{평지} & \makecell{저지} & \makecell{완경사} & \makecell{급경사} & \makecell{고지} \\
\hline
\makecell{비교치} & \makecell{1.05} & \makecell{1.00} & \makecell{0.95} & \makecell{0.90} & \makecell{0.85} \\
\hline
\end{tabular}\par}

\textbf{<자료 6> 신축단가 등} \par
1. 철근콘크리트조 슬래브지붕(부대설비 포함) @1,400,000원/m\textsuperscript{2}(2013. 1. 1) \par
(2013년 이후 건축비 변동은 없는 것으로 함) \par
2. 물리적 내용연수: 철근콘크리트조 슬래브 50년 \par
3. 잔존가치는 없는 것으로 봄 \par
\textbf{<자료 7> 기타} \par
A시 B구 근린생활시설의 최근 낙찰가율은 75\% 정도임 \par
\end{longquestion}
